\documentclass{article}
\usepackage[left=2cm, right=2cm, top=2cm, bottom=2cm]{geometry}
\usepackage{graphicx}
\usepackage{amsmath}
\usepackage{amssymb}
\usepackage{amsthm}
\usepackage{fancyhdr}
\usepackage{verbatim}
\usepackage{listings}
\usepackage{xcolor}
\usepackage{pdfpages}
\usepackage{cancel}
\usepackage{physics}
\usepackage{esint}
\usepackage{braket}

\lstset{
    language=C++,
    basicstyle=\ttfamily\footnotesize,
    keywordstyle=\color{blue},
    commentstyle=\color{green},
    stringstyle=\color{red},
    numbers=left,
    numberstyle=\tiny\color{gray},
    frame=single,
    breaklines=true,
    captionpos=b,
}


\title{IQUIST Seminar notes: 9/1}
\author{Cliff Sun}

\newtheorem{theorem}{Theorem}[section]
\newtheorem{lemma}[theorem]{Lemma}
\newtheorem{definition}[theorem]{Definition}
\newtheorem{conjecture}[theorem]{Conjecture}
\newtheorem{proposition}[theorem]{Proposition}
\newtheorem{corollary}[theorem]{Corollary}
\newtheorem{one minute paper}[theorem]{One Minute Paper}

\pagestyle{fancy}
\lhead{\textbf{\thepage}\ \ \nouppercase{\rightmark}}
\chead{IQUIST Seminar notes: Amorphous Defects and Superconducting Control of scaled quantum systems}
\rhead{Cliff Sun}

\begin{document}

\maketitle

% \section*{Lecture Span}
% \begin{enumerate}
%     \item Amorphous Defects and Superconducting Control of scaled quantum systems
% \end{enumerate}

\section*{Notes}

Hamiltonian:

\begin{equation}
    4E_c(n - n_g)^2 - E_j\cos(\pi \Phi_a)\cos(2\pi \Phi) + E_L(\Phi - \Phi_D)^2
\end{equation}

This is a flux qubit, you can also apply  

\begin{enumerate}
    \item $\alpha$ bias by tilting potential
    \item $\Delta$ bias by adding a double potential
\end{enumerate}

Questions:
\begin{enumerate}
    \item Why does $\alpha$ bias matter? (i.e, why does tilting the potential well matter?)
\end{enumerate}

You can use the $\alpha$ biasing to construct an "AND" gate. 

Room temperature and CMOS control requires many cable connections to Millikelvin qubits. They developed a control chip used for this AND gate. 

Dual rail basis is 

\begin{equation}
    \ket{0} = \ket{10}, \quad \ket{1} = \ket{01}
\end{equation}

The state is where the photon is on the "rail". 

The $CZ$ and the preparation of the zero state is due to amorphous defects in the flux qubits. They used FEM to prove that these amorphous defects are dominant. 

They found that BOE preclean and EBL at 4.7x reduced area helps with reducing the defect. They used RQL-based superconducting qubit architecture. 
\end{document}